\begin{table}[t]
\centering
\caption{五个工业基准上的零样本异常定位结果。所有指标均以百分比报告；最佳和次优结果分别用粗体和下划线标出。外部方法按其公开零样本报告口径整理，平均值按该方法可用数据集计算。}
\label{tab:main-results}
\resizebox{\textwidth}{!}{%
\begin{tabular}{lcccccccc}
\toprule
\multicolumn{9}{c}{\textbf{图像级结果 (Image AUROC, Image AP)}} \\
\midrule
数据集 & CLIP & WinCLIP & VAND & CoOp & AdaCLIP & AnomalyCLIP & AA-CLIP$^{\dagger}$ & \method{} \\
\midrule
MVTec AD & (74.1, 87.6) & (\underline{91.8}, \underline{96.5}) & (86.1, 93.5) & (88.8, 94.8) & (89.2, 96.4) & (91.6, 96.4) & (90.5, 94.9) & (\textbf{94.5}, \textbf{97.6}) \\
VisA & (66.4, 71.5) & (78.1, 81.2) & (78.0, 81.4) & (62.8, 68.1) & (\textbf{85.8}, 84.9) & (82.1, \underline{85.4}) & (\underline{84.6}, 82.2) & (\underline{84.6}, \textbf{87.4}) \\
MPDD & (54.3, 65.4) & (63.6, 69.9) & (73.0, 80.2) & (55.1, 64.2) & (76.0, 80.4) & (\underline{77.0}, \underline{82.0}) & (75.1, 80.1) & (\textbf{77.8}, \textbf{82.3}) \\
BTAD & (34.5, 52.5) & (68.2, 70.9) & (73.6, 68.6) & (66.8, 77.4) & (88.6, 92.4) & (88.3, 87.3) & (\textbf{94.8}, \textbf{97.9}) & (\underline{93.9}, \underline{94.9}) \\
DTD-Synthetic & (71.6, 85.7) & (93.2, 92.6) & (86.4, 95.0) & (-, -) & (\underline{95.5}, 97.0) & (93.5, 97.0) & (93.3, \underline{97.8}) & (\textbf{96.9}, \textbf{98.7}) \\
\midrule
\textbf{平均} & (60.2, 72.5) & (79.0, 82.2) & (79.4, 83.7) & (68.4, 76.1) & (87.0, 90.2) & (86.5, 89.6) & (\underline{87.7}, \underline{90.6}) & (\textbf{89.5}, \textbf{92.2}) \\
\midrule
\multicolumn{9}{c}{\textbf{像素级结果 (Pixel AUROC, Pixel AUPRO)}} \\
\midrule
MVTec AD & (38.4, 11.3) & (85.1, 64.6) & (87.6, 44.0) & (33.3, 6.7) & (88.7, 37.8) & (91.2, 83.2) & (\textbf{91.9}, \underline{84.6}) & (\underline{91.8}, \textbf{86.2}) \\
VisA & (46.6, 14.8) & (79.6, 56.8) & (94.2, 86.8) & (24.2, 3.8) & (\underline{95.5}, 56.8) & (\underline{95.5}, \underline{87.0}) & (\underline{95.5}, 83.0) & (\textbf{96.2}, \textbf{91.7}) \\
MPDD & (62.1, 33.0) & (76.4, 48.9) & (94.1, 83.2) & (15.4, 2.3) & (96.1, 60.3) & (96.5, \underline{88.7}) & (\underline{96.7}, 76.5) & (\textbf{97.3}, \textbf{89.9}) \\
BTAD & (30.6, 4.4) & (72.7, 27.3) & (60.8, 25.0) & (28.6, 3.8) & (92.1, 32.5) & (94.2, \underline{74.8}) & (\textbf{97.0}, 69.0) & (\underline{96.3}, \textbf{79.5}) \\
DTD-Synthetic & (33.9, 12.5) & (83.9, 57.8) & (95.3, 86.9) & (-, -) & (\underline{97.7}, 75.0) & (\textbf{97.9}, \textbf{92.3}) & (96.4, 85.9) & (\textbf{97.9}, \underline{91.8}) \\
\midrule
\textbf{平均} & (42.3, 15.2) & (79.5, 51.1) & (86.4, 65.2) & (25.4, 4.2) & (94.0, 52.5) & (95.1, \underline{85.2}) & (\underline{95.5}, 79.8) & (\textbf{95.9}, \textbf{87.8}) \\
\bottomrule
\end{tabular}%
}
\vspace{0.25em}
\footnotesize $^{\dagger}$ 表示训练式 anomaly-aware CLIP 强基线；本文对模块作用的同协议受控分析见表~\ref{tab:controlled-core-ablation}。
\end{table}
